\section{Conclusiones}
\label{sec:conclusiones}

La práctica final ha permitido ampliar y mejorar la práctica seis, que integraba las cinco prácticas previas
en un sistema multiagente cohesionado, demostrando que las técnicas
aisladas estudiadas durante el curso (predicción temporal,
clasificación de texto, OCR, agentes LLM y biometría) pueden
colaborar bajo una orquestación con LangGraph, sirviéndose
mutuamente a través de un bus de microservicios REST y bajo una
capa de observabilidad que cierra el ciclo de vida del
sistema en producción.

Los resultados cuantitativos mestran el esfuerzo invertido en las
tres evoluciones incrementales: la \textit{accuracy} cold-start del
clasificador de categoría sube del 52\,\% al 84\,\% gracias a los
\textit{embeddings} multilingües; la \textit{accuracy} del OCR
sobre facturas españolas pasa del 6{,}67\,\% al 100\,\% tras
reentrenar con datos europeos y enriquecer el contrato con
\textit{metadata} estructurada; y la biometría con vídeo más
anti-spoofing neutraliza el FAR del 36\,\% obtenido con foto única
contra ataques de pantalla. Las 21 pruebas unitarias y las pruebas
de carga con Locust garantizan que estas mejoras no han degradado
la fiabilidad ni el comportamiento bajo concurrencia.

Más allá de los números, el principal aprendizaje ha sido el de
\textbf{construir un sistema que se puede observar y operar}: el
\texttt{MonitorAgent} y el chat de administrador hacen visible el
estado interno del sistema, Langfuse permite auditar cada llamada al
LLM, y los grafos LangGraph se pueden inspeccionar dinámicamente
desde la propia UI. Esta capacidad de introspección es la que
distingue un prototipo académico de un sistema susceptible de ser
desplegado.

Como líneas de trabajo futuro destacan: (i) la finalización de las
fases 4--6 de la biometría (giros de cabeza, Moiré, rPPG); (ii) la
adopción de Alembic para gestionar las migraciones de esquema más
allá de la migración ligera implementada en \texttt{init\_db.py};
(iii) la sustitución de los \textit{tokens} en \texttt{localStorage}
por \textit{cookies} \texttt{httpOnly} con CSRF; (iv) la extracción
de los módulos \texttt{/modules/p*} a contenedores independientes en
Kubernetes para escalar el OCR de forma horizontal; y (v) el
entrenamiento de un modelo de \textit{intent classification}
específico que precalifique la decisión del orquestador y reduzca el
número de llamadas a LLM en interacciones puramente conversacionales.
